\chapter*{Resum}
\addcontentsline{toc}{chapter}{Resum}

Els mercats d'actius digitals mouen un volum econòmic considerable i, malgrat això, moltes de les seues dinàmiques encara es gestionen amb eines parcials, processos manuals o sistemes d'anàlisi poc integrats. Aquesta limitació és especialment visible en mercats amb alta fricció, on les decisions de compra, venda o manteniment depenen de senyals canviants, costos de transacció, liquiditat limitada i diferències entre plataformes.

Aquest treball proposa l'estudi d'una arquitectura multiagent basada en aprenentatge per reforç per donar suport a la presa de decisions en aquest tipus de mercats. La idea principal és distribuir el problema entre diversos agents especialitzats, de manera que cadascun puga aprendre a partir d'una part concreta de l'entorn i contribuir a una decisió global sobre la cartera. Amb aquest enfocament es busca reduir la complexitat del problema, limitar la dependència de decisions manuals i analitzar si un sistema d'agents pot arribar a generar estratègies més estables que una gestió purament humana o estàtica.

La metodologia s'aplica a l'economia virtual de \textit{Counter-Strike 2}, un entorn digital amb actius intercambiables, historial de preus, liquiditat variable i activitat freqüent en diferents mercats. El projecte combina adquisició de dades, modelatge de cartera, simulació d'operacions i entrenament d'agents amb l'objectiu d'avaluar el comportament del sistema sota condicions realistes de mercat.

\textbf{Paraules clau:} mercats d'actius digitals; presa de decisions; arquitectura multiagent; aprenentatge per reforç; gestió de cartera.

\chapter*{Resumen}
\addcontentsline{toc}{chapter}{Resumen}

Los mercados de activos digitales mueven un volumen económico considerable y, aun así, muchas de sus dinámicas siguen gestionándose mediante herramientas parciales, procesos manuales o sistemas de análisis poco integrados. Esta limitación resulta especialmente relevante en mercados con alta fricción, donde las decisiones de compra, venta o mantenimiento dependen de señales cambiantes, costes de transacción, liquidez limitada y diferencias entre plataformas.

Este trabajo plantea el estudio de una arquitectura multiagente basada en aprendizaje por refuerzo para apoyar la toma de decisiones en este tipo de mercados. La idea central consiste en distribuir el problema entre varios agentes especializados, de forma que cada uno pueda aprender a partir de una parte concreta del entorno y contribuir a una decisión global sobre la cartera. Con este enfoque se busca reducir la complejidad del problema, disminuir la dependencia de decisiones manuales y analizar si un sistema de agentes puede llegar a generar estrategias más estables que una gestión puramente humana o estática.

La metodología se aplica a la economía virtual de \textit{Counter-Strike 2}, un entorno digital con activos intercambiables, historial de precios, liquidez variable y actividad frecuente en diferentes mercados. El proyecto combina adquisición de datos, modelado de cartera, simulación de operaciones y entrenamiento de agentes con el objetivo de evaluar el comportamiento del sistema bajo condiciones realistas de mercado.

\textbf{Palabras clave:} mercados de activos digitales; toma de decisiones; arquitectura multiagente; aprendizaje por refuerzo; gestión de cartera.

\chapter*{Abstract}
\addcontentsline{toc}{chapter}{Abstract}

Digital asset markets handle a significant economic volume, yet many of their dynamics are still managed through partial tools, manual processes or poorly integrated analytical systems. This limitation is especially relevant in high-friction markets, where buy, sell or hold decisions depend on changing signals, transaction costs, limited liquidity and differences across trading platforms.

This thesis studies a multi-agent architecture based on reinforcement learning to support decision-making in this type of market. The central idea is to distribute the problem among several specialized agents, so that each one can learn from a specific part of the environment and contribute to a global portfolio decision. This approach aims to reduce the complexity of the problem, lower the dependence on manual decisions and assess whether a system of agents can produce more stable strategies than purely human or static portfolio management.

The methodology is applied to the virtual economy of \textit{Counter-Strike 2}, a digital environment with tradable assets, price histories, variable liquidity and frequent activity across different marketplaces. The project combines data acquisition, portfolio modelling, trade simulation and agent training in order to evaluate the behaviour of the system under realistic market conditions.

\textbf{Key words:} digital asset markets; decision-making; multi-agent architecture; reinforcement learning; portfolio management.
